\section{Objective 4: Synchronise Logging of Flight Telemetry and Radar Measurements}

\subsection{Introduction}

\subsection{Methodology}



\subsubsection{SBC Setup}

\begin{itemize}
    \item \textbf{Researching and selecting a suitable SBC}
    \begin{itemize}
        \item State the SBC requirements imposed by the radar application, including sufficient processing capability, memory, USB connectivity, GPIO access, size, and suitability for installation on the aircraft.
        \item Explain that a Raspberry Pi was prevalent in the reviewed literature and was therefore selected as a well-supported platform.
        \item Justify the final selection of the Raspberry Pi 5 with 16~GB of memory against the expected SDR processing requirements.
        \item Identify the supporting hardware procured for on-board use: a protective case, active cooling, and a 64~GB microSD card.
    \end{itemize}

    \item \textbf{Setting up the SBC for GPIO input and SDR interfacing}
    \begin{itemize}
        \item Describe assembling the Raspberry Pi and resolving the mechanical incompatibility between the case and the separate active cooler by using the cooler integrated into the case.
        \item Explain the selection of 64-bit Ubuntu Desktop 26.04 because the existing SDR software was intended to operate in an Ubuntu environment.
        \item Outline imaging the operating system onto the microSD card, booting the Raspberry Pi, and confirming that the installation was successful.
        \item Describe enabling remote development through SSH: installing and enabling the OpenSSH server, allowing SSH through the firewall, and connecting through the VS Code Remote--SSH extension.
        \item Explain the installation and roles of the GPIO packages (\texttt{python3-gpiozero}, \texttt{python3-lgpio}, and \texttt{gpiod}).
        \item Note that Ubuntu does not grant ordinary users GPIO access by default; document creating a \texttt{gpio} group, adding the user to it, and granting the group read/write access to \texttt{/dev/gpiochip*}.
        \item Introduce the intended physical interfaces: a GPIO input for the initial activation-signal prototype and USB for connection to the SDR.
    \end{itemize}
    \item \textbf{Writing the initial GPIO-read script}
    \begin{itemize}
        \item Define the initial purpose of the Python script: monitor a selected Raspberry Pi GPIO pin and detect the activation signal generated by the flight controller.
        \item Explain the use of rising- and falling-edge detection and \texttt{time.monotonic\_ns()} to measure PWM pulse width.
        \item State the preliminary decision thresholds: pulse widths above approximately 1700~$\mu\mathrm{s}$ represent the switch-high state and widths below approximately 1300~$\mu\mathrm{s}$ represent the switch-low state.
        \item Mention the placeholder \texttt{trigger\_response()} function and signal handlers used to clean up the GPIO state when the program exits.
        \item Record that containerisation was identified as a possible deployment step, but no implementation or verification was documented at this stage.
    \end{itemize}
\end{itemize}

\subsubsection{SBC Interface with Flight Controller and SDR}
\textit{Sub-objectives: TOBJ4-4, TOBJ4-5, TOBJ4-6, TOBJ4-7}

\begin{itemize}
    \item \textbf{Configuring Mission Planner and flight-controller logging}
    \begin{itemize}
        \item Describe installing Mission Planner, connecting the flight controller by USB, and confirming that ArduPilot firmware was already installed.
        \item Distinguish between telemetry logs stored by Mission Planner on the ground computer and DataFlash logs stored on the flight controller's SD card.
        \item Explain the observed logging behaviour when the vehicle is armed and disarmed, particularly whether GPS data are included.
        \item Discuss the \texttt{LOG\_DISARMED} parameter and why logging before arming may be important for capturing the complete synchronisation event.
        \item Outline the process for downloading a DataFlash log through MAVLink and converting or exporting it using Mission Planner.
        \item Mention the investigated Herelink network as a possible wireless route between the flight controller, SBC, and ground computer; clearly distinguish this proposed approach from the interfaces that were experimentally verified.
    \end{itemize}

    \item \textbf{Generating and verifying the activation signal}
    \begin{itemize}
        \item Explain the initial architecture: assign a transmitter switch to an RC channel, route that channel to a flight-controller PWM output, and connect the output to a Raspberry Pi GPIO input.
        \item Summarise radio calibration and transmitter channel assignment in Mission Planner, including the selection of a convenient two-position switch and the reassignment of transmitter mixes where required.
        \item Clarify the distinction between RC input channels and the numbered flight-controller servo output pins, which was an important point discovered during configuration.
        \item Describe configuring an output with the corresponding \texttt{RCIN} function so that RC channel 15 was reproduced on flight-controller output 7.
        \item State that the safety output had to be enabled before a waveform appeared, and that oscilloscope testing confirmed correct switching and a 3.3~V signal compatible with the Raspberry Pi input.
    \end{itemize}

    \item \textbf{Writing and testing the second SBC script}
    \begin{itemize}
        \item Describe updating the Python script to monitor BCM GPIO~17 and classify the measured PWM signal as switch-high or switch-low.
        \item Explain how edge callbacks measured individual pulse widths and invoked \texttt{trigger\_response()} only when the high-state threshold was detected.
        \item Note the practical networking issue encountered on the university network and the use of another network to install the required packages and update the software.
        \item Report that the flight-controller-to-Raspberry-Pi activation path was subsequently tested and detected the PWM input as expected.
    \end{itemize}

    \item \textbf{Initiating SDR data capture}
    \begin{itemize}
        \item Describe copying the existing radar application to the Raspberry Pi and first building it natively as a minimum viable implementation.
        \item List the required build dependencies (\texttt{cmake}, \texttt{pkg-config}, and \texttt{libusb-1.0-0-dev}) and report that the native build completed successfully.
        \item Explain how \texttt{trigger\_response()} was modified to launch the built \texttt{pupradar} executable through a subprocess when the activation signal was detected.
        \item State the extent of verification accurately: the Python script successfully called the radar executable, but an end-to-end recording could not yet be confirmed at that point because the SDR was not connected to the Raspberry Pi.
        \item Identify the telemetry-recording check still required: select representative fields and timestamps from an exported log and compare them with Mission Planner's live values or another known reference.
    \end{itemize}
\end{itemize}

\subsubsection{Processing Synchronisation and System Verification}
\textit{Sub-objectives: TOBJ4-8, TOBJ4-9, TOBJ4-10}

\begin{itemize}
    \item \textbf{Evaluating synchronisation methods}
    \begin{itemize}
        \item Introduce the three approaches considered: stream GPS data directly to the SBC, trigger telemetry and radar recording with the same activation signal, or align independently recorded datasets using a common elapsed-time reference.
        \item Explain the attempt to associate the selected RC switch with both radar activation and flight-controller logging through \texttt{RC15\_OPTION}, and why a suitable built-in option was not identified.
        \item Describe the alternative arming-based approach in which the Raspberry Pi records a start time and assigns each radar capture a timestamp relative to arming.
        \item Include the Lua prototype that set a servo output according to \texttt{arming:is\_armed()}, together with the required \texttt{SCR\_ENABLE} and \texttt{SERVO8\_FUNCTION} settings.
        \item Discuss why the GPIO/Lua approach was rejected: PWM output was affected by the safety switch, its relationship with GPS logging and \texttt{LOG\_DISARMED} was uncertain, and it introduced avoidable hardware and timing dependencies.
    \end{itemize}

    \item \textbf{Final MAVLink-based synchronisation architecture}
    \begin{itemize}
        \item Explain that MAVLink was installed on the Raspberry Pi and a working connection to the flight controller was established.
        \item State the architectural improvement: arming state, GPS information, and other required vehicle status can be read directly by the SBC, removing the need for the separate GPIO activation input and indirect GPS-status detection.
        \item Define the event used as the common reference time, such as receipt of the MAVLink armed-state transition, and record its timestamp using a monotonic clock.
        \item For each radar capture, record the elapsed time from the common event and retain sufficient metadata to associate the radar file with the corresponding telemetry log.
        \item Describe how telemetry timestamps and SBC-relative timestamps will be transformed to a common time axis during post-processing, including any measured or assumed offset.
        \item Specify safeguards for the final script: prevent repeated launches from one event, record connection or subprocess errors, cleanly close files, and preserve timestamps if either data stream stops unexpectedly.
    \end{itemize}

    \item \textbf{Ground integration test (TOBJ4-9)}
    \begin{itemize}
        \item Document the complete test configuration, software versions, cable/network connections, logging parameters, and the event used to start recording.
        \item While the aircraft remains on the ground, generate a sequence of identifiable events or movements visible in both datasets.
        \item Confirm that the MAVLink state change starts the intended radar-logging process and that both radar and telemetry files are created without data loss.
        \item Align the datasets using the common reference and quantify synchronisation error, repeatability, latency, and any drift over the test duration.
        \item Present plots or a table comparing matched events, then define and assess an acceptance criterion for adequate synchronisation.
    \end{itemize}

    \item \textbf{Flight integration test (TOBJ4-10)}
    \begin{itemize}
        \item State the pre-flight checks: secure the Raspberry Pi and SDR, verify cooling and power, confirm storage capacity, check all connections, and confirm telemetry and radar logging readiness.
        \item Describe the flight-test sequence and record exact arming, take-off, manoeuvre, landing, and disarming events for comparison across the two datasets.
        \item Verify that both systems record continuously throughout flight and that files remain valid after landing and shutdown.
        \item Apply the alignment method to the flight data and compare radar observations with corresponding position, altitude, attitude, and velocity telemetry.
        \item Report the achieved synchronisation accuracy, missed or duplicated captures, communication interruptions, processing limitations, and changes required for future tests.
        \item Note that the source notes for TOBJ4-9 and TOBJ4-10 contain no completed test results; populate these points only after the respective tests have been performed.
    \end{itemize}
\end{itemize}


\subsection{Results and Discussion}

\subsection{Summary}